A heterogeneidade da demanda torna a escolha da política de reposição uma decisão própria, além da calibração de seus parâmetros. Esta dissertação propõe o \textit{Adaptive Inventory Policy Engine} (AIPE), uma metodologia modular de apoio ao gerente de logística e estoques para reduzir custos evitáveis preservando o atendimento. O procedimento tem cinco etapas: especificar a decisão, representar o contexto, construir evidência comparável, recomendar uma política e avaliar suas consequências. Seu \textit{Policy Selection Engine} (PSE) associa perfis a famílias ou estima recomendações por série, com uma regra de intervenção que admite conservar a referência. A avaliação utiliza sete políticas em duas aplicações separadas. Na Bahia, 145 séries em duas origens totalizam 4.350 ajustes; a seleção por perfil aumenta o custo médio em 1,93\% e o atendimento em 3,89 pontos percentuais frente à referência global. No piloto M5, 30 séries de dez lojas em três origens totalizam 1.350 ajustes. O preditivo por classe apresenta maior custo e menor atendimento. Uma extensão exploratória por custo e serviço previstos, com recuo, reduz o custo médio em 1,56\% e aumenta o atendimento em 0,27 ponto, mas altera apenas três decisões; uma série concentra cerca de 92\% da economia. Oráculos retrospectivos distinguem erros de escolha, insuficiência do portfólio e inviabilidade física. A contribuição é um procedimento que vincula cada recomendação à evidência para subsidiar a decisão humana. A relevância prática depende do contexto; os custos avaliados são de cenário. A família permanece fixa em cada bloco de sete ciclos, sem validar revisão cíclica do PSE. Os resultados demonstram sua execução nas interfaces avaliadas e delimitam a utilidade das implementações. Coortes retrospectivas, calibrações distintas, horizontes curtos e desenvolvimento no piloto restringem a generalização e exigem confirmação independente.

\vspace{2em}
\textbf{Palavras-chave}: Metodologia de seleção; Gestão de estoques; Seleção contextual de políticas; Nível de serviço; AIPE.
